% !TeX program = pdflatex
\documentclass[12pt]{article}

% Idioma y codificación
\usepackage[spanish]{babel}
\usepackage[utf8]{inputenc}
\usepackage[T1]{fontenc}
\usepackage{lmodern}

% Página y espaciado (para acercarse a ~15 páginas)
\usepackage{geometry}
\geometry{letterpaper, margin=1in}
\usepackage{setspace}
\onehalfspacing

% Tipografía y microajustes
\usepackage{microtype}
\usepackage{csquotes}

% Tablas
\usepackage{booktabs}
\usepackage{longtable}
\usepackage{array}

% Listas
\usepackage{enumitem}

% Código
\usepackage{listingsutf8}
\usepackage{listings}
\lstset{inputencoding=utf8}
\lstset{literate=%
  {á}{{\'a}}1 {é}{{\'e}}1 {í}{{\'\i}}1 {ó}{{\'o}}1 {ú}{{\'u}}1
  {Á}{{\'A}}1 {É}{{\'E}}1 {Í}{{\'I}}1 {Ó}{{\'O}}1 {Ú}{{\'U}}1
  {ñ}{{\~n}}1 {Ñ}{{\~N}}1 {ü}{{\"u}}1 {Ü}{{\"U}}1
}
\usepackage{xcolor}
\lstdefinestyle{python}{
  language=Python,
  inputencoding=utf8,
  basicstyle=\ttfamily\small,
  numbers=left,
  numberstyle=\tiny,
  stepnumber=1,
  numbersep=8pt,
  breaklines=true,
  showstringspaces=false,
  tabsize=2,
  frame=single
}
\lstdefinestyle{sql}{
  language=SQL,
  inputencoding=utf8,
  basicstyle=\ttfamily\small,
  numbers=left,
  numberstyle=\tiny,
  stepnumber=1,
  numbersep=8pt,
  breaklines=true,
  showstringspaces=false,
  tabsize=2,
  frame=single
}

% Hipervínculos
\usepackage{hyperref}
\hypersetup{
  colorlinks=true,
  linkcolor=blue,
  urlcolor=blue,
  citecolor=blue
}

% Formato académico y paginación por sección
\usepackage{titlesec}
\usepackage{etoolbox}
\titlespacing*{\section}{0pt}{3ex plus 1ex minus .2ex}{2ex}
\titlespacing*{\subsection}{0pt}{2.5ex plus 0.8ex minus .2ex}{1.2ex}
\titlespacing*{\subsubsection}{0pt}{2ex plus 0.6ex minus .2ex}{1ex}
\preto\section{\clearpage} % cada sección inicia en nueva hoja

% Encabezados/pies limpios
\usepackage{fancyhdr}
\pagestyle{fancy}
\fancyhf{}
\lhead{Documento de apoyo (Sesión 1, Unidad 1)}
\rhead{Fundamentos de BI y OLAP}
\cfoot{\thepage}
\renewcommand{\headrulewidth}{0.4pt}
\renewcommand{\footrulewidth}{0pt}

% Bibliografía (biblatex + biber)
\usepackage[
  backend=biber,
  style=ieee,
  sorting=nyt,
  maxbibnames=99
]{biblatex}
\addbibresource{referencias_semana1.bib}

% Metadatos del documento
\title{\textbf{Documento de apoyo (Sesión 1, Unidad 1)}\\
\large Fundamentos de BI y Análisis Multidimensional\\
\normalsize Introducción a Business Intelligence y OLAP}
\author{Javier Andrés Valdez González / Auxiliar \\\\
\small Seminario de Sistemas 2 — Sección P \\
\small Universidad de San Carlos de Guatemala}
\date{28 de enero de 2026}

\begin{document}
\maketitle

\begin{abstract}
Business Intelligence (BI) reúne procesos, tecnologías y buenas prácticas para convertir datos en información útil, generar análisis y apoyar decisiones. En este contexto, OLAP (Online Analytical Processing) es una tecnología para análisis multidimensional orientada a explorar indicadores desde múltiples perspectivas (por ejemplo: tiempo, región y producto). Este documento introduce los fundamentos de BI y OLAP, distingue tipos de información por estructura y propósito, resume niveles de madurez analítica (descriptivo, diagnóstico, predictivo y prescriptivo), y explica el rol de la calidad, trazabilidad y gobernanza de datos en la toma de decisiones. Finalmente, se describen dimensiones, medidas y operaciones OLAP (drill-down, roll-up, slice, dice y pivot), así como las ventajas del análisis multidimensional para el autoservicio y la consistencia de KPI.
\end{abstract}

\textbf{Palabras clave:} Business Intelligence, OLAP, análisis multidimensional, Data Warehouse, KPI, DIKW, madurez analítica.

\newpage
\tableofcontents

% =========================================================
\section{Panorama general: BI + OLAP}
Business Intelligence (BI) reúne procesos, tecnologías y buenas prácticas para convertir datos en información útil, generar análisis y apoyar decisiones. OLAP (Online Analytical Processing) es una de las tecnologías clásicas para el análisis multidimensional, muy usada para explorar indicadores por diferentes perspectivas (por ejemplo: tiempo, región, producto).

\subsection*{Objetivos de la unidad}
\begin{itemize}[leftmargin=*, itemsep=0.25em]
  \item Reconocer qué es BI y cómo se integra con el análisis de datos.
  \item Distinguir tipos de información (estructura, propósito y nivel de detalle).
  \item Identificar niveles de madurez analítica y el valor de cada uno.
  \item Comprender cómo los datos apoyan la toma de decisiones y el control del desempeño.
  \item Introducir OLAP y el análisis multidimensional (dimensiones y medidas).
\end{itemize}

\subsection*{Arquitectura conceptual (alto nivel)}
A nivel conceptual, una solución BI suele integrar: (i) fuentes de datos operacionales y externas; (ii) un proceso de integración (ETL/ELT) para limpieza y reglas de negocio; (iii) un repositorio analítico (por ejemplo, \textit{Data Warehouse} o \textit{lakehouse}) que preserva histórico; (iv) una capa de modelado semántico (esquemas dimensionales y/o cubos); y (v) mecanismos de consumo como reportes, tableros y consultas \textit{ad-hoc}.

\begin{quote}\small
Nota: En implementaciones modernas, el DWH puede convivir con un \textit{data lake} o un \textit{lakehouse} y con herramientas de autoservicio.
\end{quote}

\noindent\textbf{Fuentes:} \cite{gartner_abi, ibm_bi, oracle_bi, gartner_bi_platforms}.

% =========================================================
\section{¿Qué es Business Intelligence?}
Business Intelligence (BI) se refiere al conjunto de capacidades que permiten acceder, integrar y analizar información para mejorar el desempeño y optimizar decisiones. BI no es una sola herramienta: incluye integración de datos, almacenamiento analítico, modelado semántico, métricas (KPI) y formas de entrega (reportes, tableros, analítica).

\subsection*{Elementos típicos de una solución BI}
\begin{itemize}[leftmargin=*, itemsep=0.25em]
  \item \textbf{Datos:} fuentes internas y externas (transaccionales y no estructuradas).
  \item \textbf{Integración:} ETL/ELT, limpieza, reglas de negocio, catalogación.
  \item \textbf{Almacenamiento:} Data Warehouse, data marts o lakehouse, con historia.
  \item \textbf{Modelos analíticos:} esquemas dimensionales, cubos OLAP, capa semántica.
  \item \textbf{Consumo:} reportes, dashboards, consultas \textit{ad-hoc} y analítica avanzada.
\end{itemize}

\subsection*{Beneficios}
Decisiones más rápidas, consistencia de indicadores, mayor visibilidad del negocio, detección de oportunidades y control de riesgos.

\noindent\textbf{Fuentes:} \cite{gartner_abi, ibm_bi, oracle_bi}.

% =========================================================
\section{Tipos de información}
Para BI es clave entender que no todos los datos son iguales. Existen diferencias por su estructura, origen y propósito (operacional vs.\ analítico). Además, BI suele trabajar con información agregada y contextualizada (no solo registros aislados).

\subsection*{Por estructura}
\begin{itemize}[leftmargin=*, itemsep=0.25em]
  \item \textbf{Estructurada:} tablas con filas/columnas (por ejemplo, ventas, inventario).
  \item \textbf{Semi-estructurada:} datos con etiquetas o marcas (JSON, XML, \textit{logs}).
  \item \textbf{No estructurada:} texto libre, imágenes, audio, correos, documentos.
\end{itemize}

\subsection*{De datos a conocimiento (DIKW)}
Un marco común es la pirámide DIKW: los datos se transforman en información cuando se organizan y adquieren significado; luego se convierten en conocimiento al interpretarse con experiencia; y apoyan sabiduría/decisión al guiar acciones.

\begin{table}[h!]
\centering
\begin{tabular}{@{}p{3cm}p{10cm}@{}}
\toprule
\textbf{Nivel} & \textbf{Enfoque} \\
\midrule
Datos & Hechos crudos, registros aislados, sin interpretación. \\
Información & Datos organizados con contexto y significado. \\
Conocimiento & Interpretación de la información (patrones, aprendizajes). \\
Sabiduría/Decisión & Aplicación del conocimiento para decidir y actuar. \\
\bottomrule
\end{tabular}
\caption{Resumen del marco DIKW.}
\end{table}

\noindent\textbf{Fuentes:} \cite{ontotext_dikw, matillion_structured}.

% =========================================================
\section{Niveles de madurez analítica}
La madurez analítica describe qué tan bien una organización usa los datos para tomar decisiones. Una forma práctica de verlo es por el tipo de análisis que domina en el día a día.

\begin{table}[h!]
\centering
\begin{tabular}{@{}p{3cm}p{4cm}p{6cm}@{}}
\toprule
\textbf{Nivel} & \textbf{Pregunta típica} & \textbf{Resultado} \\
\midrule
Descriptivo & ¿Qué pasó? & Reportes, tableros, KPI históricos \\
Diagnóstico & ¿Por qué pasó? & Análisis de causas, segmentación, correlaciones \\
Predictivo & ¿Qué podría pasar? & Modelos de pronóstico, tendencias, riesgo \\
Prescriptivo & ¿Qué deberíamos hacer? & Recomendaciones, optimización, simulaciones \\
\bottomrule
\end{tabular}
\caption{Madurez analítica por tipo de análisis dominante.}
\end{table}

Una organización más madura no ``abandona'' los niveles anteriores: los integra. La clave es pasar de reportar (describir) a anticipar y recomendar (predictivo/prescriptivo), apoyándose en datos confiables y en habilidades analíticas.

\noindent\textbf{Fuentes:} \cite{gartner_data_analytics, davenport_competing}.

% =========================================================
\section{Datos para la toma de decisiones}
En BI, los datos son un activo. Para que realmente apoyen decisiones, deben ser confiables, consistentes y disponibles a tiempo. Muchas organizaciones consolidan información en un Data Warehouse (DWH) o repositorio analítico, que integra datos de múltiples sistemas y conserva histórico.

\subsection*{Principios clave}
\begin{itemize}[leftmargin=*, itemsep=0.25em]
  \item \textbf{Calidad:} exactitud, completitud, consistencia y oportunidad.
  \item \textbf{Trazabilidad:} saber de dónde viene el dato y cómo se transformó.
  \item \textbf{Gobernanza:} roles, reglas, catálogo y definiciones estándar de KPI.
  \item \textbf{Contexto:} el significado importa (dimensiones como tiempo, producto, cliente).
\end{itemize}

\subsection*{Data Warehouse: definiciones de referencia}
Bill Inmon caracteriza un \textit{Data Warehouse} como una colección orientada a temas, integrada, variante en el tiempo y no volátil, creada para apoyar la toma de decisiones de la gerencia. En paralelo, el enfoque de Kimball enfatiza el modelado dimensional (hechos y dimensiones) para facilitar el análisis y la comprensión del negocio.

\noindent\textbf{Fuentes:} \cite{inmon_dataversity, kimball_dwtoolkit}.

% =========================================================
\section{Introducción a OLAP}
OLAP (Online Analytical Processing) es una tecnología para realizar consultas complejas y análisis multidimensional de alto desempeño sobre grandes volúmenes de datos. Su idea central es representar los datos como un ``cubo'' con dimensiones (perspectivas de análisis) y medidas (métricas numéricas).

\subsection*{Cubo OLAP: ejemplo de dimensiones}
\begin{itemize}[leftmargin=*, itemsep=0.25em]
  \item \textbf{Tiempo:} año $\rightarrow$ trimestre $\rightarrow$ mes $\rightarrow$ día
  \item \textbf{Región:} país $\rightarrow$ departamento $\rightarrow$ ciudad
  \item \textbf{Producto:} categoría $\rightarrow$ subcategoría $\rightarrow$ SKU
  \item \textbf{Medidas típicas:} ventas, unidades, margen, costo
\end{itemize}

\subsection*{Tipos comunes de OLAP}
\begin{itemize}[leftmargin=*, itemsep=0.25em]
  \item \textbf{MOLAP:} almacena estructuras multidimensionales pre-calculadas (alto rendimiento).
  \item \textbf{ROLAP:} opera sobre bases relacionales (flexible; depende del motor SQL).
  \item \textbf{HOLAP:} combina ambos enfoques.
\end{itemize}

\noindent\textbf{Fuentes:} \cite{ibm_olap, gartner_bi_platforms, kimball_dwtoolkit}.

% =========================================================
\section{Modelo multidimensional y operaciones OLAP}
El modelo multidimensional organiza la información como \textbf{hechos} (por ejemplo: ventas) y \textbf{dimensiones} (tiempo, tienda, producto, cliente). Las dimensiones suelen tener jerarquías (Año $\rightarrow$ Trimestre $\rightarrow$ Mes $\rightarrow$ Día), que permiten navegar el nivel de detalle.

\subsection*{Operaciones OLAP más usadas}
\begin{itemize}[leftmargin=*, itemsep=0.25em]
  \item \textbf{Drill-down:} ir a mayor detalle (de año a mes, de país a ciudad).
  \item \textbf{Roll-up:} resumir/agregar (de día a mes).
  \item \textbf{Slice:} fijar un valor en una dimensión (solo 2025).
  \item \textbf{Dice:} filtrar un subconjunto (2025 y Región = Centro).
  \item \textbf{Pivot:} rotar la vista (cambiar filas/columnas entre dimensiones).
\end{itemize}

\subsection*{OLTP vs.\ OLAP (comparativa)}
\begin{table}[h!]
\centering
\begin{tabular}{@{}p{6.2cm}p{6.2cm}@{}}
\toprule
\textbf{OLTP (transaccional)} & \textbf{OLAP (analítico)} \\
\midrule
Optimizado para escribir y mantener operaciones del día a día & Optimizado para leer, explorar y comparar \\
Muchas transacciones pequeñas & Histórico y agregaciones; múltiples perspectivas \\
Datos actuales, detalle fino & Datos integrados; visión consolidada por dimensiones \\
Consultas frecuentes y simples & Consultas menos frecuentes pero más complejas \\
\bottomrule
\end{tabular}
\caption{Diferencias generales entre sistemas OLTP y OLAP.}
\end{table}

\noindent\textbf{Fuentes:} \cite{ibm_olap, g4g_olap_ops}.

% =========================================================
\section{Ventajas del análisis multidimensional}
El análisis multidimensional aporta una forma natural de explorar el negocio: permite comparar métricas desde distintas perspectivas y responder preguntas iterativamente, sin necesitar un reporte nuevo por cada duda.

\subsection*{Ventajas principales}
\begin{itemize}[leftmargin=*, itemsep=0.25em]
  \item Rapidez en consultas por pre-agregaciones y estructuras optimizadas.
  \item Navegación por jerarquías (\textit{drill-down}/\textit{roll-up}) para cambiar granularidad.
  \item Consistencia de KPI cuando hay definiciones únicas en la capa semántica.
  \item Autoservicio: analistas pueden explorar datos sin depender siempre de TI.
  \item Mejor comprensión: dimensiones reflejan el lenguaje del negocio (cliente, producto, etc.).
\end{itemize}

\subsection*{Ejemplo de exploración}
Ventas por mes $\rightarrow$ filtrar Región = Occidente (\textit{slice}) $\rightarrow$ comparar productos (\textit{pivot}) $\rightarrow$ bajar a día (\textit{drill-down}) para detectar patrones.

\noindent\textbf{Fuentes:} \cite{ibm_olap, kimball_dwtoolkit}.

% =========================================================
\section{Cierre y bibliografía}
BI y OLAP permiten transformar datos en análisis que sostienen decisiones. La calidad y gobernanza de los datos son tan importantes como las herramientas. El análisis multidimensional ayuda a entender el negocio de forma exploratoria y consistente.

\subsection*{Checklist rápido para una iniciativa BI}
\begin{itemize}[leftmargin=*, itemsep=0.25em]
  \item Definir preguntas de negocio y KPI antes de diseñar reportes.
  \item Asegurar calidad, catálogo y definición estándar de métricas.
  \item Modelar dimensiones y jerarquías pensando en el usuario final.
  \item Automatizar ETL/ELT y monitorear cargas, linaje y cambios.
  \item Habilitar autoservicio con gobernanza (seguridad y roles).
\end{itemize}

\nocite{*}
\printbibliography

\end{document}
